\chapter{Limits of Functions and Continuity}

The limits of sequences studied so far let $n$ run through the natural
numbers. This closing chapter transfers the idea to functions of a
continuous variable: what value does $f(x)$ approach as $x$ approaches a
point $a$, or as $x$ grows beyond all bounds? The answer leads directly
to \emph{continuity}, the property that makes a function
graph drawable without lifting the pen, and to the Intermediate Value
Theorem, the tool that guarantees roots and market equilibria exist and
tells us how to hunt them down numerically.

The treatment stays computational: limits are evaluated with the same
limit rules as for sequences, and the indeterminate form $0/0$ is resolved
by algebra (factoring, expanding, or multiplying by a conjugate)
never by differentiation, which only enters in \volref{derivatives}. Economic
examples run through the chapter: price schedules with quantity discounts
are natural piecewise functions, and their discontinuities have a concrete
monetary meaning.

\section{Core Rules}

\begin{rulebox}{Limit of a function; one-sided limits}
$\displaystyle\lim_{x\to a}f(x)=L$ means: as $x$ gets
arbitrarily close to $a$ (from either side, without ever equalling $a$) the values $f(x)$ get arbitrarily close to $L$. The point $a$ itself
need not lie in the domain of $f$, and $f(a)$, if it exists, plays no role.

The \textbf{one-sided limits}
\[
  \lim_{x\to a^{-}}f(x)
  \qquad\text{and}\qquad
  \lim_{x\to a^{+}}f(x)
\]
approach $a$ from the left ($x<a$) and from the right ($x>a$) only.
\textbf{The two-sided limit exists exactly when both one-sided limits
exist and are equal}, and then all three coincide. For piecewise-defined
functions, always compute the two one-sided limits separately.
\end{rulebox}

\begin{rulebox}{Limit rules}
If $\lim_{x\to a}f(x)=L$ and $\lim_{x\to a}g(x)=M$ (both finite), then
\begin{itemize}[leftmargin=1.6em,itemsep=0.3ex]
  \item $\lim\limits_{x\to a}\bigl(f(x)\pm g(x)\bigr)=L\pm M$,
  \item $\lim\limits_{x\to a}\bigl(f(x)\cdot g(x)\bigr)=L\cdot M$,
        \quad in particular $\lim\limits_{x\to a}c\,f(x)=cL$,
  \item $\lim\limits_{x\to a}\dfrac{f(x)}{g(x)}=\dfrac{L}{M}$ \quad
        provided $M\neq 0$,
  \item $\lim\limits_{x\to a}f(x)^{k}=L^{k}$ and
        $\lim\limits_{x\to a}\sqrt[k]{f(x)}=\sqrt[k]{L}$ (where defined).
\end{itemize}
These are the same rules as for sequences, with $x\to a$ in place of
$n\to\infty$; they also hold for $x\to\pm\infty$ and for one-sided limits.
For all the standard functions below, limits at points of the domain are
found by \textbf{direct substitution}: $\lim_{x\to a}f(x)=f(a)$.
\end{rulebox}

\begin{rulebox}{The indeterminate form $0/0$}
If numerator and denominator both tend to $0$, the quotient rule does not
apply: the limit can still be anything. Two algebraic escapes:
\begin{itemize}[leftmargin=1.6em,itemsep=0.3ex]
  \item \textbf{Factor and cancel:} if a polynomial vanishes at $x=a$,
        then $(x-a)$ is one of its factors. Factor numerator and
        denominator, cancel the common factor $(x-a)$, then substitute.
        Example:
        \[
          \lim_{x\to 2}\frac{x^{2}-4}{x-2}
          =\lim_{x\to 2}\frac{(x-2)(x+2)}{x-2}
          =\lim_{x\to 2}\,(x+2)=4 .
        \]
  \item \textbf{Conjugate trick for roots:} expand the fraction with the
        conjugate $\sqrt{u}+v$ of a root expression $\sqrt{u}-v$ and use
        $(\sqrt{u}-v)(\sqrt{u}+v)=u-v^{2}$. Example:
        \[
          \lim_{x\to 0}\frac{\sqrt{x+9}-3}{x}
          =\lim_{x\to 0}\frac{x}{x\bigl(\sqrt{x+9}+3\bigr)}
          =\lim_{x\to 0}\frac{1}{\sqrt{x+9}+3}=\frac{1}{6} .
        \]
\end{itemize}
\end{rulebox}

\begin{rulebox}{Limits at $\pm\infty$; horizontal and vertical asymptotes}
For a rational function $\dfrac{p(x)}{q(x)}$ compare the degrees: divide
numerator and denominator by the highest power of $x$ in the denominator.
\begin{itemize}[leftmargin=1.6em,itemsep=0.3ex]
  \item $\deg p<\deg q$: \quad limit $0$ as $x\to\pm\infty$,
  \item $\deg p=\deg q$: \quad limit $=$ ratio of the leading
        coefficients,
  \item $\deg p>\deg q$: \quad the limit is $+\infty$ or $-\infty$
        (sign from the leading terms).
\end{itemize}
The line $y=L$ is a \textbf{horizontal asymptote} if $f(x)\to L$ as $x\to\infty$ or $x\to-\infty$. The line
$x=a$ is a \textbf{vertical asymptote} if at
least one one-sided limit at $a$ is $+\infty$ or $-\infty$; for rational
functions this happens where the denominator vanishes but the numerator
does not \emph{after cancelling common factors}. A common factor that
cancels completely produces a hole in the graph, not an asymptote.
\end{rulebox}

\begin{rulebox}{Sandwich rule}
If $g(x)\le f(x)\le h(x)$ for all $x$ near $a$ and
\[
  \lim_{x\to a}g(x)=\lim_{x\to a}h(x)=L,
  \qquad\text{then}\qquad
  \lim_{x\to a}f(x)=L .
\]
The same holds for $x\to\pm\infty$. The standard application is
``bounded times vanishing'': since $-1\le\sin t\le 1$ for all $t$,
\[
  -x^{2}\;\le\; x^{2}\sin\frac{1}{x}\;\le\; x^{2}
  \quad\Longrightarrow\quad
  \lim_{x\to 0}x^{2}\sin\frac{1}{x}=0 .
\]
\end{rulebox}

\begin{rulebox}{Continuity}
$f$ is \textbf{continuous at $a$} if three things hold:
$f(a)$ is defined, $\lim_{x\to a}f(x)$ exists, and
\[
  \lim_{x\to a}f(x)=f(a) .
\]
$f$ is \textbf{continuous on an interval} if it is continuous at every
point of the interval. All the standard functions are continuous wherever
they are defined: polynomials and rational functions, roots $\sqrt[k]{x}$,
powers $x^{r}$, the absolute value $|x|$, exponentials $\mathrm{e}^{x}$
and $a^{x}$, the logarithm $\ln x$, and $\sin x$, $\cos x$. Moreover,
\textbf{sums, differences, products, quotients (where the denominator is
nonzero), and compositions of continuous functions are continuous}. In
practice, discontinuities can therefore only occur at points where a
formula breaks down (denominator $=0$, log of $0$, \dots) or where a
piecewise definition switches branches.
\end{rulebox}

\begin{rulebox}{Classification of discontinuities}
Suppose $f$ is not continuous at $a$. Compute the one-sided limits.
\begin{itemize}[leftmargin=1.6em,itemsep=0.3ex]
  \item \textbf{Removable}: the
        two-sided limit $L$ exists (finite), but $f(a)\neq L$ or $f(a)$
        is undefined. Redefining $f(a):=L$ repairs the function:
        the graph has a single hole. Typical source: a cancelled $0/0$.
  \item \textbf{Jump}: both one-sided limits exist
        and are finite but differ. The jump has size
        $\lim_{x\to a^{+}}f(x)-\lim_{x\to a^{-}}f(x)$.
  \item \textbf{Pole}: at least one one-sided limit is
        $+\infty$ or $-\infty$; the line $x=a$ is a vertical asymptote.
\end{itemize}
Piecewise price schedules are the standard economic example: if a
quantity discount applies to \emph{all} units once the order passes a
threshold, the total bill jumps \emph{downward} at the threshold: a
jump discontinuity that makes some order sizes irrational to choose. If
the discount applies only to the units \emph{above} the threshold, the
bill is continuous.
\end{rulebox}

\begin{rulebox}{Intermediate Value Theorem (IVT)}
Let $f$ be continuous on $[a,b]$. Then $f$
attains every value between $f(a)$ and $f(b)$. In particular:
\[
  f(a)\cdot f(b)<0
  \quad\Longrightarrow\quad
  f(x^{*})=0 \text{ for some } x^{*}\in(a,b).
\]
This is the existence tool for roots and for market equilibria (apply it
to excess demand $D(p)-S(p)$). It also drives the \textbf{bisection
method}: evaluate $f$ at the midpoint $m=\tfrac{a+b}{2}$ and keep the
half-interval on which $f$ still changes sign. Each step halves the
bracket, so after $n$ steps the root is located to within
$(b-a)/2^{n}$. Continuity is essential: a function with a jump can change
sign without ever being zero.
\end{rulebox}

\begin{rulebox}{Weierstrass: extreme values on $[a,b]$}
A function continuous on a \emph{closed, bounded} interval $[a,b]$
attains a maximum and a minimum there.
This is the one-variable case of the extreme value theorem that later
volumes drill in several variables; here it guarantees, before any
optimisation technique is available, that ``best'' and ``worst'' values
exist at all. Both hypotheses matter: $f(x)=1/x$ on $(0,1]$ has no
maximum, and $f(x)=x$ on $[0,\infty)$ has none either.
\end{rulebox}

\section{Worked Example}

\begin{worked}{Evaluate $\lim_{x\to0}\dfrac{\sqrt{x+4}-2}{x}$, choose $a$ and $b$ making a piecewise function continuous, and classify the discontinuities of three examples}
\wstep{Step 1. Substitute first --- and read what the substitution says.}
At $x=0$ the expression is $\frac{2-2}{0}=\frac00$. That is \textbf{not} an
answer and \textbf{not} a proof that the limit fails to exist; it is a
signal that numerator and denominator share a vanishing factor, which
algebra must expose.

\wstep{Step 2. Rationalise the numerator.}
A difference of square roots is cleared by multiplying by its conjugate:
\[
  \frac{\sqrt{x+4}-2}{x}\cdot\frac{\sqrt{x+4}+2}{\sqrt{x+4}+2}
  =\frac{(x+4)-4}{x\bigl(\sqrt{x+4}+2\bigr)}
  =\frac{x}{x\bigl(\sqrt{x+4}+2\bigr)} .
\]
The factor $x$ has been made visible --- which was the entire purpose.
\wtrap{Multiply the \emph{numerator} by its conjugate when the roots are on
top. The reflex of rationalising the denominator is unhelpful here; the
goal is not tidiness but exposing the common factor.}

\wstep{Step 3. Cancel, then substitute.}
For $x\neq0$ --- and $x\neq0$ is exactly the situation a limit concerns ---
\[
  \frac{x}{x(\sqrt{x+4}+2)}=\frac{1}{\sqrt{x+4}+2}
  \xrightarrow[x\to0]{}\frac{1}{2+2}=\boxed{\tfrac14}.
\]
Numerical check: at $x=0.01$ the original quotient is $0.2499\ldots$
\checkmark

\wstep{Step 4. Now make $f$ continuous at $x=0$.}
Define
\[
  f(x)=\begin{cases}\dfrac{\sqrt{x+4}-2}{x}, & x\neq0,\\[1.2ex] a, & x=0.\end{cases}
\]
Continuity at a point requires three things: $f(0)$ exists, the limit
exists, and they are \emph{equal}. The first two hold, so the condition is
$a=\tfrac14$. This is a \textbf{removable} discontinuity: one value, chosen
correctly, repairs it.

\wstep{Step 5. A two-sided version, where both a limit and a slope must match.}
\[
  g(x)=\begin{cases}x^{2}, & x\le1,\\ ax+b, & x>1.\end{cases}
\]
Continuity at $1$ needs the one-sided limits to agree:
\[
  \lim_{x\to1^{-}}g=1,\qquad \lim_{x\to1^{+}}g=a+b
  \quad\Longrightarrow\quad a+b=1 .
\]
If differentiability is also demanded, the slopes must match too:
$2x|_{x=1}=2$ from the left and $a$ from the right, so $a=2$ and hence
$b=-1$.
\wtrap{Continuity gives \emph{one} equation, differentiability a second.
With two unknowns you need both conditions --- and note the order: a
function can be continuous without being differentiable ($|x|$ at $0$), but
never differentiable without being continuous.}

\wstep{Step 6. Classify three discontinuities.}
\begin{center}\footnotesize
\begin{tabular}{lll}
  \hline
  function & at & type\\
  \hline
  $\dfrac{x^{2}-1}{x-1}$ & $x=1$ & removable: both one-sided limits are $2$\\[0.8ex]
  $\dfrac{1}{x-2}$ & $x=2$ & infinite (pole): $\to-\infty$ left, $+\infty$ right\\[0.8ex]
  $\sgn(x)$ & $x=0$ & jump: limits $-1$ and $+1$, both finite, unequal\\
  \hline
\end{tabular}
\end{center}
\wtrap{All three have ``a problem at a point'', and they behave completely
differently. Decide by computing the two one-sided limits: equal and finite
$\Rightarrow$ removable; finite but different $\Rightarrow$ jump; at least
one infinite $\Rightarrow$ pole. Only the removable kind can be patched.}

\wstep{Step 7. Continuity earns two theorems --- use them.}
\emph{Intermediate Value Theorem.} $h(x)=x^{3}+x-1$ is continuous, with
$h(0)=-1<0$ and $h(1)=1>0$. So $h$ takes the value $0$ somewhere in
$(0,1)$: the equation has a root, proved without solving it. This is how
every existence claim about an equilibrium price is made.

\emph{Weierstrass.} A continuous function on a \textbf{closed and bounded}
interval attains a maximum and a minimum. Both hypotheses are load-bearing:
$1/x$ on $(0,1]$ is continuous but unbounded, because the interval is not
closed; $x$ on $[0,\infty)$ has no maximum, because the interval is not
bounded.
\wtrap{In \volref{optimization} the existence half of every optimisation problem is
Weierstrass, and the usual failure is a feasible set that is open ($x>0$
instead of $x\ge0$) or unbounded. When an optimisation problem ``has no
solution'', check those two words before checking the calculus.}
\end{worked}

\section{Practice Problems}

Evaluate each limit or answer the question posed. Where a limit fails to
exist, say so and give the one-sided limits; write $+\infty$ or $-\infty$
where appropriate.

\begin{problems}
  \item $\displaystyle\lim_{x\to 2}\frac{x^{2}+1}{x+3}$
  \item $\displaystyle\lim_{x\to 2}\frac{x^{2}-4}{x-2}$
  \item $\displaystyle\lim_{x\to -3}\frac{x^{2}+5x+6}{x+3}$
  \item $\displaystyle\lim_{x\to 1}\frac{x^{3}-1}{x-1}$
  \item $\displaystyle\lim_{x\to 2}\frac{x^{2}-5x+6}{x^{2}-4}$
  \item $\displaystyle\lim_{x\to 4}\frac{\sqrt{x}-2}{x-4}$
  \item $\displaystyle\lim_{x\to 0}\frac{\sqrt{x+9}-3}{x}$
  \item For $f(x)=\dfrac{|x-2|}{x-2}$ compute
        $\displaystyle\lim_{x\to 2^{-}}f(x)$ and
        $\displaystyle\lim_{x\to 2^{+}}f(x)$. Does
        $\displaystyle\lim_{x\to 2}f(x)$ exist?
  \item For
        \[
          f(x)=
          \begin{cases}
            x^{2}, & x<1,\\
            3-x,   & x\ge 1,
          \end{cases}
        \]
        compute both one-sided limits at $x=1$. Does the two-sided limit
        exist? Classify the discontinuity, if any.
  \item $\displaystyle\lim_{x\to\infty}\frac{3x^{2}-5x+1}{2x^{2}+7}$
        \quad and \quad
        $\displaystyle\lim_{x\to-\infty}\frac{3x^{2}-5x+1}{2x^{2}+7}$
  \item $\displaystyle\lim_{x\to\infty}\frac{2x+1}{x^{2}-4}$. Which
        horizontal asymptote does this give?
  \item $\displaystyle\lim_{x\to\infty}\frac{x^{3}-2x}{5x^{2}+1}$
        \quad and \quad
        $\displaystyle\lim_{x\to-\infty}\frac{x^{3}-2x}{5x^{2}+1}$
  \item Find all horizontal and vertical asymptotes of
        $f(x)=\dfrac{2x^{2}-8}{x^{2}-1}$, including the one-sided
        behaviour at each vertical asymptote.
  \item Find all asymptotes of $f(x)=\dfrac{x^{2}-9}{x-3}$ and classify
        the discontinuity at $x=3$.
  \item Use the sandwich rule to compute
        $\displaystyle\lim_{x\to 0}\,x^{2}\sin\frac{1}{x}$.
  \item Find the constant $c$ that makes
        \[
          f(x)=
          \begin{cases}
            cx+3,    & x<2,\\
            x^{2}+1, & x\ge 2,
          \end{cases}
        \]
        continuous at $x=2$ (and hence on all of $\R$).
  \item Classify all discontinuities of
        $f(x)=\dfrac{x+2}{x^{2}-4}$.
  \item A supplier offers a quantity discount on the \emph{whole} order:
        \[
          C(q)=
          \begin{cases}
            12q, & 0\le q\le 50,\\
            9q,  & q>50,
          \end{cases}
        \]
        where $C(q)$ is the total bill for $q$ units.
        (a)~Compute the one-sided limits of $C$ at $q=50$ and classify
        the discontinuity.
        (b)~The supplier switches to $C(q)=9q+F$ for $q>50$ with a fixed
        fee $F$. Which $F$ makes $C$ continuous at $q=50$, and what does
        the continuous schedule mean economically?
  \item Show with the Intermediate Value Theorem that
        $f(x)=x^{3}-3x+1$ has a root in $[0,1]$, then perform one
        bisection step: evaluate $f$ at the midpoint and state the new,
        halved bracket.
  \item On which set is
        $g(x)=\dfrac{\sqrt{x+2}}{x-1}$
        continuous? Justify using the continuity of standard functions
        and of compositions and quotients.
\end{problems}

\section{Solutions}

\begin{solutions}
  \item The denominator does not vanish at $x=2$, so substitute
        directly:
        \[
          \lim_{x\to 2}\frac{x^{2}+1}{x+3}=\frac{4+1}{2+3}=1 .
        \]

  \item At $x=2$ both numerator and denominator are $0$: the form is
        $0/0$. Factor and cancel:
        \[
          \frac{x^{2}-4}{x-2}=\frac{(x-2)(x+2)}{x-2}=x+2
          \quad (x\neq 2),
          \qquad
          \lim_{x\to 2}\,(x+2)=4 .
        \]

  \item Form $0/0$. Since $x^{2}+5x+6=(x+2)(x+3)$,
        \[
          \lim_{x\to -3}\frac{(x+2)(x+3)}{x+3}
          =\lim_{x\to -3}\,(x+2)=-1 .
        \]

  \item Form $0/0$. Use $x^{3}-1=(x-1)(x^{2}+x+1)$:
        \[
          \lim_{x\to 1}\frac{(x-1)(x^{2}+x+1)}{x-1}
          =\lim_{x\to 1}\,(x^{2}+x+1)=1+1+1=3 .
        \]

  \item Form $0/0$; both quadratics contain the factor $(x-2)$:
        \[
          \frac{x^{2}-5x+6}{x^{2}-4}
          =\frac{(x-2)(x-3)}{(x-2)(x+2)}
          =\frac{x-3}{x+2}
          \quad (x\neq 2),
        \]
        so the limit is $\dfrac{2-3}{2+2}=-\dfrac{1}{4}$.

  \item Form $0/0$. Expand with the conjugate $\sqrt{x}+2$:
        \[
          \frac{\sqrt{x}-2}{x-4}
          =\frac{(\sqrt{x}-2)(\sqrt{x}+2)}{(x-4)(\sqrt{x}+2)}
          =\frac{x-4}{(x-4)(\sqrt{x}+2)}
          =\frac{1}{\sqrt{x}+2},
        \]
        hence the limit is $\dfrac{1}{\sqrt{4}+2}=\dfrac{1}{4}$.

  \item Form $0/0$. Conjugate $\sqrt{x+9}+3$:
        \[
          \frac{\sqrt{x+9}-3}{x}
          =\frac{(x+9)-9}{x\bigl(\sqrt{x+9}+3\bigr)}
          =\frac{1}{\sqrt{x+9}+3}
          \;\xrightarrow[x\to 0]{}\;
          \frac{1}{3+3}=\frac{1}{6} .
        \]

  \item For $x<2$ we have $|x-2|=-(x-2)$, so $f(x)=-1$; for $x>2$ we
        have $|x-2|=x-2$, so $f(x)=+1$. Hence
        \[
          \lim_{x\to 2^{-}}f(x)=-1,
          \qquad
          \lim_{x\to 2^{+}}f(x)=+1 .
        \]
        The one-sided limits differ, so $\lim_{x\to 2}f(x)$ does not
        exist ($f$ has a jump of size $2$ at $x=2$).

  \item Left branch: $\lim_{x\to 1^{-}}x^{2}=1$. Right branch:
        $\lim_{x\to 1^{+}}(3-x)=2$, and $f(1)=3-1=2$. The one-sided
        limits are finite but different, so the two-sided limit does not
        exist: $f$ has a \textbf{jump discontinuity} at $x=1$ with jump
        $2-1=1$.

  \item Degrees are equal, so the limit is the ratio of the leading
        coefficients. Dividing by $x^{2}$:
        \[
          \frac{3x^{2}-5x+1}{2x^{2}+7}
          =\frac{3-\frac{5}{x}+\frac{1}{x^{2}}}{2+\frac{7}{x^{2}}}
          \;\longrightarrow\;\frac{3}{2}
          \qquad (x\to\infty\ \text{and}\ x\to-\infty).
        \]
        Both limits equal $\tfrac{3}{2}$.

  \item The numerator has smaller degree than the denominator, so
        \[
          \lim_{x\to\infty}\frac{2x+1}{x^{2}-4}
          =\lim_{x\to\infty}
           \frac{\frac{2}{x}+\frac{1}{x^{2}}}{1-\frac{4}{x^{2}}}
          =\frac{0}{1}=0 .
        \]
        The horizontal asymptote is $y=0$ (the $x$-axis).

  \item The numerator has larger degree. Divide by $x^{2}$:
        \[
          \frac{x^{3}-2x}{5x^{2}+1}
          =\frac{x-\frac{2}{x}}{5+\frac{1}{x^{2}}} .
        \]
        The denominator tends to $5$ while the numerator tends to
        $+\infty$ resp.\ $-\infty$, so
        \[
          \lim_{x\to\infty}\frac{x^{3}-2x}{5x^{2}+1}=+\infty,
          \qquad
          \lim_{x\to-\infty}\frac{x^{3}-2x}{5x^{2}+1}=-\infty .
        \]

  \item \textbf{Horizontal:} equal degrees, ratio of leading
        coefficients $\tfrac{2}{1}=2$, so $y=2$ as $x\to\pm\infty$.
        \textbf{Vertical:} the denominator $x^{2}-1=(x-1)(x+1)$
        vanishes at $x=\pm 1$; the numerator $2x^{2}-8$ equals $-6\neq 0$
        at both points, so no cancellation: both are poles. Near $x=1$
        the numerator is close to $-6$, and $(x-1)(x+1)$ is a small
        positive number for $x\to 1^{+}$, a small negative number for
        $x\to 1^{-}$:
        \[
          \lim_{x\to 1^{+}}f(x)=-\infty,\qquad
          \lim_{x\to 1^{-}}f(x)=+\infty .
        \]
        Near $x=-1$ the numerator is again close to $-6$, while
        $(x-1)(x+1)$ is negative for $x\to-1^{+}$ and positive for
        $x\to-1^{-}$:
        \[
          \lim_{x\to -1^{+}}f(x)=+\infty,\qquad
          \lim_{x\to -1^{-}}f(x)=-\infty .
        \]
        Asymptotes: $y=2$, $x=1$, $x=-1$.

  \item Factor: $x^{2}-9=(x-3)(x+3)$, so for $x\neq 3$
        \[
          f(x)=\frac{(x-3)(x+3)}{x-3}=x+3,
          \qquad
          \lim_{x\to 3}f(x)=6 .
        \]
        The limit exists but $f(3)$ is undefined: the discontinuity at
        $x=3$ is \textbf{removable} (a hole at $(3,6)$), \emph{not} a
        vertical asymptote, the common factor cancels completely.
        Since $f(x)=x+3\to\pm\infty$ as $x\to\pm\infty$, there is no
        horizontal asymptote either: $f$ has no asymptotes at all.

  \item Since $-1\le\sin\frac{1}{x}\le 1$ for all $x\neq 0$ and
        $x^{2}\ge 0$,
        \[
          -x^{2}\;\le\; x^{2}\sin\frac{1}{x}\;\le\; x^{2}.
        \]
        Both bounds tend to $0$ as $x\to 0$, so by the sandwich rule
        \[
          \lim_{x\to 0}\,x^{2}\sin\frac{1}{x}=0 .
        \]

  \item Both branches are polynomials, hence continuous away from
        $x=2$; only the seam matters. One-sided limits at $x=2$:
        \[
          \lim_{x\to 2^{-}}(cx+3)=2c+3,
          \qquad
          \lim_{x\to 2^{+}}(x^{2}+1)=5=f(2).
        \]
        Continuity requires $2c+3=5$, i.e.\ $c=1$. With $c=1$
        both one-sided limits equal $f(2)=5$, so $f$ is continuous on
        $\R$.

  \item Factor the denominator: $x^{2}-4=(x-2)(x+2)$, so for
        $x\neq -2$
        \[
          f(x)=\frac{x+2}{(x-2)(x+2)}=\frac{1}{x-2}.
        \]
        \textbf{At $x=-2$:} the factor cancels and
        $\lim_{x\to-2}f(x)=\frac{1}{-4}=-\frac14$ exists, while $f(-2)$
        is undefined: a \textbf{removable} discontinuity.
        \textbf{At $x=2$:} after cancelling, the denominator still
        vanishes while the numerator is $1$:
        \[
          \lim_{x\to 2^{-}}\frac{1}{x-2}=-\infty,
          \qquad
          \lim_{x\to 2^{+}}\frac{1}{x-2}=+\infty,
        \]
        a \textbf{pole} (vertical asymptote $x=2$).

  \item (a) One-sided limits at the threshold:
        \[
          \lim_{q\to 50^{-}}C(q)=12\cdot 50=600=C(50),
          \qquad
          \lim_{q\to 50^{+}}C(q)=9\cdot 50=450 .
        \]
        Both are finite but different: a \textbf{jump discontinuity}
        with jump $450-600=-150$. The bill \emph{drops} by $150$ as the
        order passes $50$ units: ordering $51$ units costs
        $9\cdot 51=459$, less than the $600$ charged for $50$ units,
        indeed less than the $12\cdot 39=468$ charged for $39$ units.
        No rational customer orders between about $39$ and $50$ units.

        (b) Continuity at $q=50$ requires
        \[
          \lim_{q\to 50^{+}}(9q+F)=450+F\overset{!}{=}600
          \quad\Longrightarrow\quad
          F=150 .
        \]
        Then for $q>50$,
        \[
          C(q)=9q+150=12\cdot 50+9\,(q-50):
        \]
        the first $50$ units cost $12$ each and only the units
        \emph{beyond} the threshold get the discounted price $9$. The
        continuous schedule is exactly the marginal-discount tariff.

  \item $f(x)=x^{3}-3x+1$ is a polynomial, hence continuous on $[0,1]$.
        Endpoint values:
        \[
          f(0)=1>0,
          \qquad
          f(1)=1-3+1=-1<0 .
        \]
        The sign changes, so by the IVT there is an $x^{*}\in(0,1)$ with
        $f(x^{*})=0$. \textbf{Bisection step:} midpoint $m=\tfrac12$,
        \[
          f\Bigl(\tfrac12\Bigr)
          =\tfrac18-\tfrac32+1=-\tfrac38<0 .
        \]
        The sign change now sits between $f(0)>0$ and
        $f(\tfrac12)<0$: the new bracket is $\bigl[0,\tfrac12\bigr]$,
        half as long as before.

  \item The numerator $\sqrt{x+2}$ is the composition of the continuous
        functions $x\mapsto x+2$ and $t\mapsto\sqrt{t}$, hence
        continuous wherever it is defined, i.e.\ for $x+2\ge 0$. The
        denominator $x-1$ is continuous everywhere, and a quotient of
        continuous functions is continuous wherever the denominator is
        nonzero. Therefore $g$ is continuous on its whole domain
        \[
          [-2,1)\cup(1,\infty),
        \]
        i.e.\ for all $x\ge -2$ with $x\neq 1$.
\end{solutions}

\section{Additional Problems}

\begin{problems}
  \item $\displaystyle\lim_{x\to -2}\frac{x^{2}-x-6}{x+2}$
  \item $\displaystyle\lim_{x\to 9}\frac{\sqrt{x}-3}{x-9}$
  \item $\displaystyle\lim_{x\to 0}\frac{\sqrt{4+x}-2}{x}$
  \item $\displaystyle\lim_{x\to\infty}\frac{5x^{3}+2x}{2x^{3}-x^{2}+4}$
        \quad and \quad
        $\displaystyle\lim_{x\to-\infty}\frac{5x^{3}+2x}{2x^{3}-x^{2}+4}$
  \item Find all asymptotes of $f(x)=\dfrac{3x+1}{x-2}$, including the
        one-sided behaviour at the vertical asymptote.
  \item Find the constant $c$ that makes
        \[
          f(x)=
          \begin{cases}
            cx^{2}, & x\le 2,\\
            3x+2,   & x>2,
          \end{cases}
        \]
        continuous on $\R$.
  \item Classify all discontinuities of
        $f(x)=\dfrac{x^{2}-1}{x^{2}-3x+2}$.
  \item Use the sandwich rule to compute
        $\displaystyle\lim_{x\to\infty}\frac{2+\cos x}{x}$.
  \item Show that $f(x)=x^{3}+x-3$ has a root in $[1,2]$ and perform
        one bisection step, stating the new bracket.
  \item Explain why $f(x)=x^{2}-4x+1$ attains a maximum and a minimum
        on $[0,3]$, and find both values (no derivatives: complete the
        square).
\end{problems}

\section{Additional Solutions}

\begin{solutions}
  \item Form $0/0$. Factor $x^{2}-x-6=(x-3)(x+2)$:
        \[
          \lim_{x\to -2}\frac{(x-3)(x+2)}{x+2}
          =\lim_{x\to -2}\,(x-3)=-5 .
        \]

  \item Form $0/0$. Conjugate $\sqrt{x}+3$, using
        $x-9=(\sqrt{x}-3)(\sqrt{x}+3)$:
        \[
          \frac{\sqrt{x}-3}{x-9}=\frac{1}{\sqrt{x}+3}
          \;\xrightarrow[x\to 9]{}\;
          \frac{1}{3+3}=\frac{1}{6} .
        \]

  \item Form $0/0$. Conjugate $\sqrt{4+x}+2$:
        \[
          \frac{\sqrt{4+x}-2}{x}
          =\frac{(4+x)-4}{x\bigl(\sqrt{4+x}+2\bigr)}
          =\frac{1}{\sqrt{4+x}+2}
          \;\xrightarrow[x\to 0]{}\;
          \frac{1}{4} .
        \]

  \item Equal degrees: the limit is the ratio of the leading
        coefficients,
        \[
          \lim_{x\to\pm\infty}\frac{5x^{3}+2x}{2x^{3}-x^{2}+4}
          =\frac{5}{2}
        \]
        in both directions (divide numerator and denominator by
        $x^{3}$).

  \item \textbf{Horizontal:} equal degrees give
        $y=\tfrac{3}{1}=3$ as $x\to\pm\infty$. \textbf{Vertical:} the
        denominator vanishes at $x=2$, the numerator equals $7\neq 0$
        there, so $x=2$ is a pole. Near $x=2$ the numerator is close to
        $7>0$, so the sign follows $x-2$:
        \[
          \lim_{x\to 2^{-}}f(x)=-\infty,
          \qquad
          \lim_{x\to 2^{+}}f(x)=+\infty .
        \]
        Asymptotes: $y=3$ and $x=2$.

  \item Both branches are continuous; match them at the seam $x=2$:
        \[
          \lim_{x\to 2^{-}}cx^{2}=4c=f(2),
          \qquad
          \lim_{x\to 2^{+}}(3x+2)=8 .
        \]
        Continuity requires $4c=8$, so $c=2$.

  \item Factor: $x^{2}-1=(x-1)(x+1)$ and $x^{2}-3x+2=(x-1)(x-2)$, so
        for $x\neq 1$
        \[
          f(x)=\frac{(x-1)(x+1)}{(x-1)(x-2)}=\frac{x+1}{x-2}.
        \]
        \textbf{At $x=1$:} the factor cancels,
        $\lim_{x\to 1}f(x)=\frac{2}{-1}=-2$ exists while $f(1)$ is
        undefined: \textbf{removable}. \textbf{At $x=2$:} the numerator
        tends to $3\neq 0$ while the denominator vanishes:
        \[
          \lim_{x\to 2^{-}}\frac{x+1}{x-2}=-\infty,
          \qquad
          \lim_{x\to 2^{+}}\frac{x+1}{x-2}=+\infty,
        \]
        a \textbf{pole} with vertical asymptote $x=2$.

  \item Since $-1\le\cos x\le 1$, we have $1\le 2+\cos x\le 3$, and for
        $x>0$
        \[
          \frac{1}{x}\;\le\;\frac{2+\cos x}{x}\;\le\;\frac{3}{x} .
        \]
        Both bounds tend to $0$ as $x\to\infty$, so the limit is $0$ by
        the sandwich rule.

  \item $f(x)=x^{3}+x-3$ is continuous (polynomial), and
        \[
          f(1)=1+1-3=-1<0,
          \qquad
          f(2)=8+2-3=7>0,
        \]
        so by the IVT a root lies in $(1,2)$. \textbf{Bisection step:}
        midpoint $m=1.5$,
        \[
          f(1.5)=3.375+1.5-3=1.875>0 .
        \]
        The sign change is now between $f(1)<0$ and $f(1.5)>0$: new
        bracket $[1,\,1.5]$.

  \item $f$ is a polynomial, hence continuous, and $[0,3]$ is a closed
        bounded interval, so by the Weierstrass theorem $f$ attains a
        maximum and a minimum on $[0,3]$. Completing the square,
        \[
          f(x)=x^{2}-4x+1=(x-2)^{2}-3 ,
        \]
        which is smallest when $(x-2)^{2}=0$: minimum $-3$ at $x=2$
        (which lies in $[0,3]$). The maximum of $(x-2)^{2}$ over
        $[0,3]$ occurs at the endpoint farther from $2$, namely $x=0$:
        $f(0)=1$ versus $f(3)=9-12+1=-2$. Maximum $1$ at $x=0$,
        minimum $-3$ at $x=2$.
\end{solutions}
